\textbf{COLUNI 2026 -- Questão 15}\par\medskip Leia os textos abaixo.

\textbf{TEXTO I}

Qual é o preço que o planeta paga por cada palavra gerada pela Inteligência Artificial?

Gerar um texto de 100 palavras no ChatGPT consome, em média, 519 mililitros de água. O consumo é porque os servidores que operam em centros de dados geram enormes quantidades de calor à medida que efetuam os milhares de cálculos necessários para cada resposta. Para evitar o sobreaquecimento, o sistema requer um arrefecimento constante, muitas vezes através de sistemas que utilizam água para transferir o calor para torres de arrefecimento.

\noindent PARRA, Sérgio. \textbf{ChatGPT: a quantidade de água consumida pela IA é alarmante}. National Geographic Portugal, 2025. Acesso em: 28 jul 2025. Disponível em: https://www.nationalgeographic.pt/meio-ambiente/sede-chatgpt-quantidade-agua-consumida-pela-ia-e-alarmante\_5618.

\textbf{TEXTO II}

Segundo o Sistema Nacional de Informações sobre Saneamento, do Ministério das Cidades, cada brasileiro consome, em média, 154 litros de água todos os dias.

\noindent \textbf{Brasileiro consome, em média, 154 litros de água por dia, aponta ONU}. Confederação Nacional de Municípios, 2018. Acesso em: 28 jul 2025. Disponível em: https://cnm.org.br/comunicacao/noticias/brasileiro-consome-em-media-154-litros-de-agua-por-dia-aponta-onu.

Suponha que 160 estudantes do CAp-Coluni fizeram, cada um, uma pesquisa de 2000 palavras no ChatGPT. A quantidade de água necessária para o arrefecimento decorrente dessa pesquisa poderia ser utilizada, em 1 dia, por no máximo:
\begin{enumerate}
\item 9 brasileiros.
\item 10 brasileiros.
\item 11 brasileiros.
\item 12 brasileiros.
\end{enumerate}
